\chapter{System Design}
\label{ch:design}

This chapter describes the design of RetailPulse as implemented. All figures in it were redrawn for this report from the current code; the figures of the mid report, which the mid evaluation found incorrect, are not reused (Appendix~\ref{app:feedback} lists what was wrong with each).

\section{Architecture}
RetailPulse uses a layered architecture (Figure~\ref{fig:architecture}). Each layer depends only on the layer below it, and every layer above the data layer reaches the database through a defined interface: the vision layer writes through the ingestion contract, the agents read through five read-only tools, and dispatch writes through one handler.

\begin{figure}[htbp]
\centering
\begin{tikzpicture}[font=\small, scale=0.92, transform shape,
  comp/.style={box, fill=white, text width=#1},
  comp/.default=3.5cm,
  lay/.style={draw, rounded corners=4pt, fill=#1, inner sep=4pt},
  dl/.style={font=\scriptsize, fill=white, inner sep=1pt}]
% presentation
\node[comp=5.2cm] (mgrui) at (3.4,0) {Manager workspace\\\scriptsize dashboard, analytics, alerts, settings, status};
\node[comp=5.2cm] (empui) at (10.6,0) {Employee workspace\\\scriptsize own tasks, shift, profile};
\node[comp=13.2cm] (proxy) at (7.0,-1.45) {Session and API proxy\\\scriptsize signed HTTP-only session, role check, adds the internal service secret};
% application
\node[comp=3.9cm] (dash) at (2.2,-3.75) {Dashboard and Ask routes\\\scriptsize scoped to the manager's store};
\node[comp=3.9cm] (api) at (7.0,-3.75) {Backend API\\\scriptsize accounts, alerts, settings, tasks};
\node[comp=3.9cm] (jobs) at (11.8,-3.75) {Camera and CV-job routes\\\scriptsize at most 2 jobs at once};
% agents and operations
\node[comp=3.0cm] (orch) at (1.9,-6.1) {Orchestrator and business agents};
\node[comp=3.0cm] (signal) at (5.6,-6.1) {Signal policy and dispatch};
\node[comp=6.6cm] (tools) at (3.75,-7.75) {Five read-only tools\\\scriptsize footfall, dwell, interactions, unattended, coverage};
% vision
\node[comp=3.9cm] (det) at (11.8,-6.05) {Detection and tracking\\\scriptsize YOLOv8n, BoT-SORT};
\node[comp=3.9cm] (role) at (11.8,-7.45) {Role and interaction\\\scriptsize classifier, proximity, pose};
\node[comp=3.9cm] (ing) at (11.8,-8.85) {Ingestion\\\scriptsize events, health, heatmaps};
% data
\node[box, double, double distance=1.2pt, text width=12.2cm, fill=white] (db) at (7.0,-11.0) {\textbf{PostgreSQL}\\\scriptsize application tables: accounts, rosters, zones, alerts, tasks, dispatch decisions\\\scriptsize event tables: tracks, interactions, unattended episodes, camera health, heatmaps};
% layers
\begin{scope}[on background layer]
\node[lay=blue!6, fit=(mgrui)(empui)(proxy), label={[font=\scriptsize\bfseries, anchor=south west, inner sep=1pt]north west:Presentation (Next.js)}] {};
\node[lay=green!6, fit=(dash)(api)(jobs), label={[font=\scriptsize\bfseries, anchor=south west, inner sep=1pt]north west:Application (FastAPI)}] {};
\node[lay=orange!8, fit=(orch)(signal)(tools), label={[font=\scriptsize\bfseries, anchor=south west, inner sep=1pt]north west:Agents and operations}] {};
\node[lay=violet!6, fit=(det)(role)(ing), label={[font=\scriptsize\bfseries, anchor=south west, inner sep=1pt]north west:Vision pipeline}] {};
\node[lay=gray!10, fit=(db), label={[font=\scriptsize\bfseries, anchor=south west, inner sep=1pt]north west:Data}] {};
\end{scope}
% flows
\draw[flow] (mgrui) -- (mgrui |- proxy.north);
\draw[flow] (empui) -- (empui |- proxy.north);
\foreach \t in {dash, api, jobs} { \draw[flow] (\t |- proxy.south) -- (\t); }
\draw[flow] (dash.south -| orch) -- (orch);
\draw[flow] (orch) -- (orch |- tools.north);
\draw[flow] (signal) -- node[dl, right] {reads} (signal |- tools.north);
\draw[flow] (jobs) -- (det);
\draw[flow] (det) -- (role);
\draw[flow] (role) -- (ing);
\draw[flow] (tools) -- node[dl, right] {read only} (tools |- db.north);
\draw[flow] (signal.east) -- ++(0.45,0) coordinate (sw) -- (sw |- db.north);
\draw[flow] ([xshift=15mm]api.south) coordinate (aw) -- (aw |- db.north);
\draw[flow] (ing) -- node[dl, right] {writes events} (ing |- db.north);
\end{tikzpicture}
\caption[Layered architecture of RetailPulse]{Layered architecture of RetailPulse. Arrows point from caller to callee. Only the backend API (accounts, alerts, tasks), dispatch (tasks, decisions, alerts) and ingestion (events) write to PostgreSQL; the agents only read.}
\label{fig:architecture}
\end{figure}

\begin{table}[H]
\centering
\caption{Architectural components and responsibilities}\label{tab:components}
\begin{tabular}{P{0.25\textwidth}P{0.42\textwidth}P{0.25\textwidth}}
\toprule
Component & Responsibility & Main source\\
\midrule
Session and API proxy & Signs sessions, checks role, forwards requests with the internal secret and actor identity & \path{frontend/src/app/api/}\\
Backend API & Accounts, invitations, alerts, settings, rosters, tasks and the task event stream & \path{backend/app.py}, \path{store_core.py}\\
Interaction pipeline & Detects, tracks and classifies people; nominates and judges interactions & \path{backend/interaction/}\\
Ingestion & Converts a run's outputs into event rows, camera health and heatmaps & \path{backend/events/ingest.py}\\
Read-only tools & Five typed, scoped analytics readers over the event store & \path{backend/agents/tools/}\\
Orchestrator & Routes a question, runs one business agent, checks every claim & \path{backend/agents/graph.py}\\
Business agents & Customer Behavior and Employee Activity specialists; digest; occupancy & \path{backend/agents/business/}\\
Signal policy & Decides when a window becomes an unattended or unstaffed event & \path{signal_policy.py}\\
Dispatch & Selects an eligible employee, creates the task, prevents duplicates & \path{backend/dispatch.py}\\
\bottomrule
\end{tabular}
\end{table}

\subsection{Agents in the Architecture}
The proposal named seven agent responsibilities. Table~\ref{tab:agent-map} shows where each one is realised. Two are full business agents; the others are realised as deterministic components, because rules are easier to test and explain than a language model wherever the decision is a fixed policy.

\begin{table}[H]
\centering
\caption{Agent responsibilities and their realisation}\label{tab:agent-map}
\begin{tabular}{P{0.27\textwidth}P{0.65\textwidth}}
\toprule
Responsibility & Realisation\\
\midrule
Customer Behavior Agent & Business specialist reading footfall, dwell, interactions and unattended episodes\\
Employee Activity Agent & Business specialist reading staff coverage, interactions and unattended episodes\\
Interaction Agent & The interaction pipeline: proximity nomination plus pose judge\\
Employee Task Assignment Agent & Signal policy plus dispatch service (deterministic eligibility rules)\\
Insight Narration Agent & Digest builder over verified claims, with a narrative number check\\
Query Agent & The orchestrator's question path over the two specialists\\
Orchestrator & LangGraph state graph: route, specialist, ground, repair, finish\\
\bottomrule
\end{tabular}
\end{table}

\section{Vision Pipeline Design}
For every processed frame the pipeline runs the steps in Figure~\ref{fig:pipeline}. A COCO-trained YOLOv8n model detects people and BoT-SORT keeps their track identifiers. A second detector trained on the team's store footage labels each new track as customer or employee during its first second and rechecks it every two seconds, so an identity switch cannot carry an employee label onto a customer. The role decision combines that detector's vote with bib colour and standing still, using the fitted weights adopted after evaluation (Section~\ref{sec:role-weights}).

\begin{figure}[htbp]
\centering
\begin{tikzpicture}[font=\small, node distance=5mm and 7mm,
  pstep/.style={box, text width=2.6cm, fill=white},
  dec/.style={draw, diamond, aspect=2.2, align=center, font=\scriptsize, inner sep=1pt}]
\node[pstep] (in) {Frame from camera or recording};
\node[pstep, right=of in] (detect) {Detect people\\\scriptsize YOLOv8n (COCO)};
\node[pstep, right=of detect] (track) {Track\\\scriptsize BoT-SORT with re-ID};
\node[pstep, right=of track] (roles) {Classify role\\\scriptsize role detector, bib, stillness};
\node[pstep, below=of roles] (prox) {Nominate pairs\\\scriptsize distance in body heights, $\geq 2$\,s};
\node[dec, left=of prox] (cand) {Employee and\\customer close?};
\node[pstep, below=10mm of cand] (judge) {Pose check\\\scriptsize head within $60^\circ$ in $\geq$ half of 4 frames};
\node[pstep, left=of judge, xshift=-6mm] (rej) {Verdict: rejected\\\scriptsize re-judged every 1\,s};
\node[pstep, right=of judge, xshift=6mm] (conf) {Verdict: interaction};
\node[pstep, below=of judge, text width=4.4cm] (out) {Write frame record, annotated frame and camera health};
\draw[flow] (in) -- (detect);
\draw[flow] (detect) -- (track);
\draw[flow] (track) -- (roles);
\draw[flow] (roles) -- (prox);
\draw[flow] (prox) -- (cand);
\draw[flow] (cand) -- node[lbl, right] {yes} (judge);
\draw[flow] (cand.west) -- node[lbl, above] {no} (cand.west -| in.south) -- ([yshift=-5mm]in.south |- out.south) -| (out.south);
\draw[flow] (judge) -- node[lbl, above] {not facing} (rej);
\draw[flow] (judge) -- node[lbl, above] {facing} (conf);
\draw[flow] (rej) |- (out);
\draw[flow] (conf) |- (out);
\end{tikzpicture}
\caption{Per-frame processing in the interaction pipeline}
\label{fig:pipeline}
\end{figure}

Distance is measured in body heights rather than metres: the allowed gap is 1.0 body height between feet, 0.8 between hips and 0.6 between heads, because higher body parts carry less depth information. A pair must stay within that gap for two seconds and must not drift apart faster than walking pace, which removes most pass-bys. If the pose model finds no face or fails, the verdict stays \emph{unverified} rather than being guessed.

\section{Question Answering and Grounding}
A manager's question passes through the orchestration graph in Figure~\ref{fig:graph}. The router either selects one business agent or ends the run as unsupported; questions that mention dispatching, ranking or identifying people are refused before routing. The agent may call only its allowed tools. Every number the agent returns is a \emph{claim} that names a tool call and the path of one metric in its result; the grounding step compares each claim with that metric. A failed claim triggers one repair attempt, and anything that still fails is dropped. The graph never writes to the database.

\begin{figure}[htbp]
\centering
\begin{tikzpicture}[font=\small, node distance=6mm and 9mm,
  st/.style={box, rounded corners=6pt, fill=white, text width=2.3cm}]
\node[circle, fill=black, inner sep=3pt] (s) {};
\node[st, right=of s] (route) {route\\\scriptsize offline rules or provider};
\node[st, right=14mm of route, yshift=9mm] (cb) {Customer Behavior agent};
\node[st, right=14mm of route, yshift=-9mm] (ea) {Employee Activity agent};
\node[st, right=14mm of cb, yshift=-9mm] (ground) {ground\\\scriptsize claim = metric?};
\node[st, below=12mm of ground] (repair) {repair\\\scriptsize once};
\node[st, right=of ground] (finish) {finish\\\scriptsize answer, citations, limits};
\node[circle, draw, inner sep=3pt, right=of finish] (e) {}; \node[circle, fill=black, inner sep=1.6pt] at (e) {};
\draw[flow] (s) -- (route);
\draw[flow] (route) -- (cb);
\draw[flow] (route) -- (ea);
\draw[flow] (route.south) -- ++(0,-2.1) -| node[lbl, pos=0.25] {unsupported or refused} (finish.south);
\draw[flow] (cb) -- (ground);
\draw[flow] (ea) -- (ground);
\draw[flow] (ground) -- node[lbl, right] {issues, first time} (repair);
\draw[flow] (repair.west) -- ++(-4mm,0) |- ([yshift=-1mm]ground.west);
\draw[flow] (ground) -- node[lbl, above] {ok} (finish);
\draw[flow] (finish) -- (e);
\end{tikzpicture}
\caption{Orchestration graph for a manager question}
\label{fig:graph}
\end{figure}

% Design models: domain model, class diagrams, data design, sequence diagrams, state machine, deployment.
\tikzset{
  cls/.style={draw, rectangle split, rectangle split parts=2, rectangle split part align={center,left},
              align=left, font=\scriptsize, fill=white, text width=#1, inner sep=2.5pt,
              execute at begin node={\hyphenpenalty=10000\exhyphenpenalty=10000}},
  cls/.default=2.9cm,
  dcls/.style={draw, rectangle split, rectangle split parts=3, rectangle split part align={center,left,left},
               align=left, font=\scriptsize, fill=white, text width=#1, inner sep=2.5pt,
               execute at begin node={\hyphenpenalty=10000\exhyphenpenalty=10000}},
  dcls/.default=4.6cm,
  rel/.style={thin},
  m/.style={font=\tiny, inner sep=1pt},
  an/.style={font=\tiny\itshape, fill=white, inner sep=1pt},
  compo/.style={thin, {Diamond[fill=black,length=2.6mm,width=1.8mm]}-},
  aggr/.style={thin, {Diamond[open,length=2.6mm,width=1.8mm]}-},
  inherit/.style={thin, -{Triangle[open,length=2.4mm,width=2.4mm]}},
  realize/.style={thin, dashed, -{Triangle[open,length=2.4mm,width=2.4mm]}},
  usedep/.style={thin, dashed, ->},
}

\section{Domain Model}
The domain model (Figure~\ref{fig:domain}) shows the concepts of the problem and how they relate. It contains only real-world concepts and their attributes: software components such as agents, the orchestrator or the database do not appear, and no operations are shown.

\begin{landscape}
\begin{figure}[p]
\centering
\begin{tikzpicture}[scale=0.95, transform shape]
% row 1
\node[cls] (health) at (0,0) {\textbf{CameraHealth}\nodepart{second}minute\\frameRate\\unknownRoleShare\\frozenOrDark};
\node[cls] (camera) at (4.6,0) {\textbf{Camera}\nodepart{second}name\\role\\source};
\node[cls] (store) at (9.2,0) {\textbf{Store}\nodepart{second}name};
\node[cls] (zone) at (13.8,0) {\textbf{Zone}\nodepart{second}name\\type\\capacity\\minimumStaff};
\node[cls] (roster) at (18.4,0) {\textbf{RosterEntry}\nodepart{second}shift\\available};
% row 2
\node[cls] (run) at (4.6,-3.4) {\textbf{AnalysisRun}\nodepart{second}startedAt\\source};
\node[cls=3.5cm] (user) at (18.4,-3.4) {\textbf{StaffMember}\nodepart{second}name\\email\\role (Manager \textbar\ Employee)\\active};
% row 3
\node[cls] (inter) at (0,-6.8) {\textbf{Interaction}\nodepart{second}start\\end\\verdict};
\node[cls] (track) at (4.6,-6.8) {\textbf{Track}\nodepart{second}role\\roleScore\\firstSeen\\lastSeen};
\node[cls] (unatt) at (9.2,-6.8) {\textbf{UnattendedEpisode}\nodepart{second}waitedSeconds\\helped};
\node[cls] (event) at (13.8,-6.8) {\textbf{StaffingEvent}\nodepart{second}type\\window\\severity};
% row 4
\node[cls] (alert) at (9.2,-10.2) {\textbf{Alert}\nodepart{second}category\\severity\\status};
\node[cls] (decision) at (13.8,-10.2) {\textbf{DispatchDecision}\nodepart{second}outcome\\reason};
\node[cls] (task) at (18.4,-10.2) {\textbf{Task}\nodepart{second}instruction\\state\\acknowledgedAfter};

% associations: label in the middle, multiplicities near the ends
\draw[rel] (health) -- node[an, above] {reports on} node[m, below, pos=0.12] {0..*} node[m, below, pos=0.88] {1} (camera);
\draw[rel] (camera) -- node[an, above] {belongs to} node[m, below, pos=0.12] {1..*} node[m, below, pos=0.88] {1} (store);
\draw[rel] (store) -- node[an, above] {has} node[m, below, pos=0.12] {1} node[m, below, pos=0.88] {1..*} (zone);
\draw[rel] (zone) -- node[an, above] {staffed by} node[m, below, pos=0.12] {0..1} node[m, below, pos=0.88] {0..*} (roster);
\draw[rel] ([xshift=6mm]camera.south) -- ++(0,-0.7) -| node[an, pos=0.25] {watches} node[m, right, pos=0.92] {1} ([xshift=-6mm]zone.south)
  node[m, left] at ([xshift=6mm,yshift=-2mm]camera.south) {0..*};
\draw[rel] ([xshift=-6mm]camera.south) -- node[an, left] {produces} node[m, left, pos=0.15] {1} node[m, left, pos=0.85] {0..*} ([xshift=-6mm]run.north);
\draw[rel] (roster) -- node[an, right] {for} node[m, right, pos=0.15] {0..*} node[m, right, pos=0.85] {1} (user);
\draw[rel] (store.north) -- ++(0,0.7) -| node[an, pos=0.25] {employs} node[m, left, pos=0.97] {1..*} ([xshift=2.4cm]user.east) -- (user.east)
  node[m, right] at ([yshift=2mm]store.north) {1};
\draw[rel] (run) -- node[an, left] {contains} node[m, left, pos=0.15] {1} node[m, left, pos=0.85] {0..*} (track);
\draw[rel] ([yshift=3mm]track.west) -- node[an, above] {employee} node[m, below, pos=0.15] {1} node[m, below, pos=0.85] {0..*} ([yshift=3mm]inter.east);
\draw[rel] ([yshift=-3mm]track.west) -- node[an, below] {customer} node[m, below, pos=0.15] {1} node[m, below, pos=0.85] {0..*} ([yshift=-3mm]inter.east);
\draw[rel] (track) -- node[an, above] {waits as} node[m, below, pos=0.12] {1} node[m, below, pos=0.88] {0..1} (unatt);
\draw[rel] (unatt) -- node[an, above] {triggers} node[m, below, pos=0.12] {1..*} node[m, below, pos=0.88] {0..1} (event);
\draw[rel] ([xshift=6mm]zone.south) -- node[an, right] {concerns} node[m, right, pos=0.1] {1} node[m, right, pos=0.9] {0..*} ([xshift=6mm]event.north);
\draw[rel] (event) -- node[an, right] {decided by} node[m, right, pos=0.15] {1} node[m, right, pos=0.85] {1} (decision);
\draw[rel] (decision) -- node[an, above] {raises} node[m, below, pos=0.12] {1} node[m, below, pos=0.88] {0..1} (alert);
\draw[rel] (decision) -- node[an, above] {creates} node[m, below, pos=0.12] {1} node[m, below, pos=0.88] {0..1} (task);
\draw[rel] (task) -- node[an, right] {assigned to} node[m, right, pos=0.12] {0..*} node[m, right, pos=0.88] {1} (user);
\draw[rel] (alert.south) -- ++(0,-0.7) -| node[an, pos=0.25] {resolved by} node[m, left, pos=0.97] {0..1} ([xshift=2.9cm,yshift=-4mm]user.east) -- ([yshift=-4mm]user.east)
  node[m, right] at ([yshift=-2mm]alert.south) {0..*};
\end{tikzpicture}
\caption{Domain model}
\label{fig:domain}
\end{figure}
\end{landscape}

\section{Class Diagrams}
The class diagrams show the main design classes with their real names, typed attributes and operations with return types. Figure~\ref{fig:class-vision} covers the vision pipeline and Figure~\ref{fig:class-agents} the agent and operations layer. Helper classes and private members are omitted.

\begin{figure}[htbp]
\centering
\begin{tikzpicture}[scale=0.9, transform shape]
% the pipeline and its configuration (left)
\node[dcls=4.3cm] (cfg) at (0,3.3) {\textbf{InteractionConfig}
  \nodepart{second}min\_interaction\_seconds: float = 2.0\\facing\_angle\_degrees: float = 60\\min\_facing\_fraction: float = 0.5\\role\_weights: dict[str, float]
  \nodepart{third}from\_json(path: Path): InteractionConfig};
\node[dcls=4.3cm] (pipe) at (0,-0.6) {\textbf{InteractionPipeline}
  \nodepart{second}config: InteractionConfig\\judge: PoseJudge \textbar\ None
  \nodepart{third}process\_frame(frame: ndarray,\\\quad timestamp: float): dict\\verdict(event: InteractionEvent): str};
% the four parts (middle column)
\node[dcls=4.5cm] (tracker) at (6.1,3.3) {\textbf{ByteTrackPersonTracker}
  \nodepart{second}tracker: str = "botsort\_reid.yaml"
  \nodepart{third}detect\_and\_track(frame: ndarray,\\\quad timestamp: float): list[dict]};
\node[dcls=4.5cm] (roles) at (6.1,1.05) {\textbf{RoleClassifier}
  \nodepart{second}agent: CustEmpClassifier
  \nodepart{third}classify(frame: ndarray, bbox: tuple,\\\quad track\_id: int, timestamp: float):\\\quad tuple[str, float, dict]};
\node[dcls=4.5cm] (eng) at (6.1,-1.45) {\textbf{EngagementDetector}
  \nodepart{second}config: InteractionConfig
  \nodepart{third}update(people: list[TrackedPerson],\\\quad timestamp: float, frame\_height: int):\\\quad list[InteractionEvent]};
\node[dcls=4.5cm] (judge) at (6.1,-3.95) {\textbf{PoseJudge}
  \nodepart{second}facing\_angle: float\\min\_facing\_fraction: float
  \nodepart{third}classify\_interaction(frames: list,\\\quad boxes: list): JudgeResult};
% results and helpers (right column)
\node[dcls=3.7cm] (cec) at (11.6,1.05) {\textbf{CustEmpClassifier}
  \nodepart{second}weights: dict[str, float]\\threshold: float
  \nodepart{third}decide(factors: dict):\\\quad tuple[str, float, dict]};
\node[dcls=3.7cm] (event) at (11.6,-1.45) {\textbf{InteractionEvent}
  \nodepart{second}employee\_track\_id: int\\customer\_track\_id: int\\started\_at: float\\distance\_body\_heights: float\\active: bool
  \nodepart{third}};
\node[dcls=3.7cm] (jr) at (11.6,-3.95) {\textbf{JudgeResult}
  \nodepart{second}label: str\\confidence: float
  \nodepart{third}};
\draw[usedep] (pipe) -- node[an, right] {configured by} (cfg);
% composition: one diamond at the whole, a bus to each part
\draw[compo] (pipe.east) -- ++(0.8,0) coordinate (bus);
\draw[thin] (bus |- tracker) -- (bus |- judge);
\draw[thin] (bus |- tracker) -- node[m, above, pos=0.8] {1} (tracker.west);
\draw[thin] (bus |- roles) -- node[m, above, pos=0.8] {1} (roles.west);
\draw[thin] (bus |- eng) -- node[m, above, pos=0.8] {1} (eng.west);
\draw[thin] (bus |- judge) -- node[m, above, pos=0.8] {0..1} (judge.west);
\draw[compo] (roles) -- node[m, above, pos=0.8] {1} (cec);
\draw[usedep] (eng) -- node[an, above] {creates} (event);
\draw[usedep] (judge) -- node[an, above] {returns} (jr);
\end{tikzpicture}
\caption{Class diagram: vision pipeline}
\label{fig:class-vision}
\end{figure}

\begin{figure}[htbp]
\centering
\begin{tikzpicture}[scale=0.9, transform shape]
% question path (left and middle)
\node[dcls=4.4cm] (svc) at (0,0) {\textbf{AnalyticsService}
  \nodepart{second}traces: MemoryTraceStore
  \nodepart{third}ask(request: AskRequest,\\\quad scope: TrustedAccessScope): GraphRun\\get\_run(run\_id: str,\\\quad scope: TrustedAccessScope): GraphRun};
\node[dcls=3.9cm] (spec) at (0,-3.7) {\textbf{SpecialistRegistration}
  \nodepart{second}handler: Callable\\allowed\_tools: frozenset[str]\\development\_stub: bool
  \nodepart{third}};
\node[dcls=4.4cm] (reg) at (5.8,0) {\textbf{ScopedToolRegistry}
  \nodepart{second}scope: TrustedAccessScope
  \nodepart{third}invoke(name: str, arguments: dict):\\\quad ToolResult\\input\_schemas(): dict};
\node[dcls=4.4cm] (ea) at (5.8,-3.7) {\textbf{EventAnalytics}
  \nodepart{second}dsn: str\\synthetic: bool
  \nodepart{third}get\_footfall(scope, request): FootfallResult\\get\_zone\_dwell(scope, request):\\\quad ZoneDwellResult\\get\_interactions(scope, request):\\\quad InteractionsResult};
\node[dcls=4.4cm] (pg) at (5.8,-7.6) {\textbf{PostgresAnalyticsTools}
  \nodepart{second}contexts: tuple[ObservationContext]
  \nodepart{third}get\_unattended(scope, request):\\\quad UnattendedResult\\get\_coverage(scope, request):\\\quad CoverageResult};
% operations path (right and bottom)
\node[dcls=4.1cm] (pol) at (11.4,0) {\textbf{SignalPolicy}
  \nodepart{second}min\_wait\_s: float = 180\\high\_wait\_s: float = 360\\min\_staffed\_ratio: float = 0.25\\max\_unknown\_share: float = 0.2
  \nodepart{third}};
\node[dcls=4.1cm] (dec) at (11.4,-3.7) {\textbf{SignalDecision}
  \nodepart{second}event\_type: str\\zone\_id: str\\reason: ReasonCode\\event: OperationalEvent \textbar\ None
  \nodepart{third}};
\node[dcls=4.1cm] (oe) at (11.4,-7.6) {\textbf{OperationalEvent}
  \nodepart{second}event\_id: str\\idempotency\_key: str\\zone\_id: str\\event\_type: str\\severity: str
  \nodepart{third}};
\node[dcls=4.1cm] (ad) at (11.4,-11.3) {\textbf{OperationalDispatchAdapter}
  \nodepart{second}handler: DispatchHandler
  \nodepart{third}handle(event: OperationalEvent,\\\quad scope: TrustedOperationalScope):\\\quad DispatchReceipt};
\node[dcls=4.4cm] (dh) at (5.8,-11.3) {\guillemotleft interface\guillemotright\\\textbf{DispatchHandler}
  \nodepart{second}
  \nodepart{third}\_\_call\_\_(event: OperationalEvent,\\\quad scope: TrustedOperationalScope):\\\quad DispatchOutcome};
\node[dcls=4.4cm] (disp) at (0,-11.3) {\guillemotleft module\guillemotright\\\textbf{dispatch}
  \nodepart{second}ACTIVE\_TASK\_STATES: tuple[str]
  \nodepart{third}handle\_event(conn, body: dict): dict\\select\_dispatch\_employee(conn,\\\quad target\_zone\_id: str): Row \textbar\ None\\make\_dispatch\_handler(dsn: str):\\\quad DispatchHandler};
\draw[usedep] (svc) -- node[an, above] {\guillemotleft create\guillemotright} (reg);
\draw[aggr] (svc) -- node[m, left, pos=0.85] {0..2} (spec);
\draw[aggr] (reg) -- node[m, right, pos=0.85] {1} (ea);
\draw[inherit] (pg) -- (ea);
\draw[usedep] (dec) -- node[an, right] {\guillemotleft use\guillemotright\ thresholds} (pol);
\draw[usedep] (dec.west) -- ++(-0.5,0) |- node[an, pos=0.75] {\guillemotleft use\guillemotright\ reads tools} (reg.east);
\draw[aggr] (dec) -- node[m, right, pos=0.85] {0..1} (oe);
\draw[usedep] (ad) -- node[an, right] {handles} (oe);
\draw[aggr] (ad) -- node[m, above, pos=0.8] {1} (dh);
\draw[usedep] (disp) -- node[an, above] {\guillemotleft create\guillemotright} (dh);
\end{tikzpicture}
\caption{Class diagram: agents and operations}
\label{fig:class-agents}
\end{figure}

\section{Data Design}
All persistent data is in one PostgreSQL database with two groups of tables. The \emph{application tables} (Figure~\ref{fig:erd}) hold accounts, store configuration and the operational record of every dispatch decision. The \emph{event tables} (Table~\ref{tab:event-tables}) hold what the cameras observed; they share one set of common columns and are linked by run, camera and track identifiers rather than foreign keys, because a run is always written and replaced as a whole.

\begin{figure}[htbp]
\centering
\begin{tikzpicture}[scale=0.88, transform shape,
  tbl/.style={draw, rectangle split, rectangle split parts=2, rectangle split part align={center,left},
              align=left, font=\scriptsize, fill=white, text width=3.3cm, inner sep=2.5pt,
              rectangle split part fill={gray!15,white}},
  fk/.style={thin, -{Stealth[length=1.6mm]}},
  m/.style={font=\tiny, inner sep=1pt}]
\node[tbl] (users) at (0,0) {\textbf{users}\nodepart{second}\underline{id}\\name, email\\role, password\_hash\\active};
\node[tbl] (rosters) at (4.6,0) {\textbf{rosters}\nodepart{second}\underline{id}\\user\_id (FK)\\zone\_id (FK, null = floater)\\shift, available};
\node[tbl] (zones) at (9.2,0) {\textbf{zones}\nodepart{second}\underline{id}\\name, type\\capacity, min\_staff};
\node[tbl] (cameras) at (13.8,0) {\textbf{cameras}\nodepart{second}\underline{id}\\zone\_id (FK)\\role, source\_type, online};
\node[tbl] (events) at (13.8,-4.6) {\textbf{orchestration\_events}\nodepart{second}\underline{id}\\event\_type, zone\_id\\severity, evidence\_json\\dedupe\_key (unique)};
\node[tbl] (dd) at (9.2,-4.6) {\textbf{dispatch\_decisions}\nodepart{second}\underline{id}\\event\_id (FK)\\alert\_id, task\_id (FK)\\selected\_employee\_id (FK users)\\decision, reason};
\node[tbl] (tasks) at (4.6,-4.6) {\textbf{tasks}\nodepart{second}\underline{id}\\employee\_id (FK)\\target\_zone\_id (FK)\\trigger\_alert\_id (FK)\\instruction, state, ack\_seconds};
\node[tbl] (tev) at (0,-4.6) {\textbf{task\_events}\nodepart{second}\underline{id}\\task\_id (FK)\\actor\_id (FK)\\state, note, created\_at};
\node[tbl] (alerts) at (9.2,-8.6) {\textbf{alerts}\nodepart{second}\underline{id}\\category, severity\\status, evidence\_json\\resolved\_by (FK users)};
% foreign keys point from the referencing table to the referenced table
\draw[fk] (rosters) -- (users);
\draw[fk] (rosters) -- (zones);
\draw[fk] (cameras) -- (zones);
\draw[fk] (tasks) -- (users);
\draw[fk] (tasks) -- (zones);
\draw[fk] (tev) -- (tasks);
\draw[fk] (tev) -- (users);
\draw[fk] (dd) -- (events);
\draw[fk] (dd) -- (tasks);
\draw[fk] (dd) -- (alerts);
\draw[fk] (tasks.south) |- ([yshift=2mm]alerts.west);
\end{tikzpicture}
\caption[Application tables]{Application tables. Arrows run from a foreign key to the table it references; the two references to \texttt{users} from \texttt{alerts} and \texttt{dispatch\_decisions} are named in the column instead of drawn. Only key columns are shown.}
\label{fig:erd}
\end{figure}

\begin{table}[H]
\centering
\caption{Event tables (common columns: event\_id, run\_id, store\_id, camera\_id, zone\_id, ts\_start, ts\_end)}\label{tab:event-tables}
\begin{tabular}{P{0.24\textwidth}P{0.68\textwidth}}
\toprule
Table & Additional content\\
\midrule
tracks & track\_id, role, employee\_score, frames\_observed\\
role\_decisions & track\_id, role, frames\\
interactions & employee and customer track\_id, duration\_s, verdict (interaction, no\_interaction, unverified)\\
unattended\_episodes & customer\_track\_id, waited\_s, helped\\
zone\_dwell / cell\_dwell & track\_id, role, dwell\_s (per zone or per registered map cell)\\
footfall & track\_id, direction (in or out)\\
queue\_samples & queue\_length, est\_wait\_s\\
pipeline\_health & fps, mean\_confidence, unknown\_role\_share, id\_churn, frozen\_or\_dark (one row per camera per minute)\\
heatmaps & cols, rows, cells (one grid per camera per hour)\\
\bottomrule
\end{tabular}
\end{table}

Two rules protect the data. Every event row must pass the contract check \texttt{validate()}: the common identifiers are non-empty and no row ends before it starts. And \texttt{dedupe\_key} on \texttt{orchestration\_events} is unique, so the same staffing event can never be inserted twice; the insert uses \texttt{ON CONFLICT DO NOTHING} and the handler then returns the earlier decision.

\section{Interaction Design}
\tikzset{
  life/.style={draw, fill=white, font=\scriptsize, align=center, minimum height=7mm, inner sep=2pt, text width=1.9cm},
  msg/.style={->, thin, >={Stealth[length=1.8mm]}},
  ret/.style={->, thin, dashed, >={Stealth[length=1.8mm]}},
  ml/.style={font=\tiny, above, inner sep=1pt, align=center},
}
\subsection{Staffing Dispatch}
Figure~\ref{fig:seq-dispatch} shows how a closed observation window becomes a task (UC-12). The operations cycle runs after each window for every configured zone. The policy reads three tool results; only if a customer waited past the threshold and the camera could see enough of the window does it emit an event. The dispatch handler then works inside one transaction: it records the event, which fails harmlessly for a duplicate, selects an eligible employee, and writes the alert, task, task event and decision together.

\begin{figure}[htbp]
\centering
\begin{tikzpicture}[scale=0.86, transform shape]
\foreach \x/\n/\t in {0/cyc/Operations\\cycle, 2.45/pol/Signal policy, 4.9/tools/Read-only\\tools, 7.35/ad/Dispatch\\adapter, 9.8/hd/Dispatch\\handler, 12.25/db/PostgreSQL, 14.7/emp/Employee\\page} {
  \node[life] (\n) at (\x,0) {\t};
  \draw[dashed, gray] (\n.south) -- (\x,-13.6);
}
\draw[msg] (0,-1.0) -- node[ml] {evaluate\_zone(zone, window)} (2.45,-1.0);
\draw[msg] (2.45,-1.6) -- node[ml] {get\_unattended, get\_zone\_dwell,\\get\_coverage} (4.9,-1.6);
\draw[msg] (4.9,-2.2) -- node[ml] {read-only SELECT} (12.25,-2.2);
\draw[ret] (12.25,-2.7) -- node[ml] {rows} (4.9,-2.7);
\draw[ret] (4.9,-3.2) -- node[ml] {typed results with evidence} (2.45,-3.2);
\node[draw, font=\tiny, fill=yellow!10, align=left, text width=2.2cm] at (1.2,-4.2) {wait $\geq$ min\_wait\_s,\\unknown share $\leq$ max};
\draw[ret] (2.45,-5.0) -- node[ml] {SignalDecision(event)} (0,-5.0);
\draw[msg] (0,-5.6) -- node[ml] {handle(event, scope)} (7.35,-5.6);
\draw[msg] (7.35,-6.2) -- node[ml] {handler(event, scope)} (9.8,-6.2);
\draw[msg] (9.8,-6.8) -- node[ml] {INSERT event ON CONFLICT (dedupe\_key) DO NOTHING} (12.25,-6.8);
% alt fragment
\draw (8.9,-7.3) rectangle (13.3,-12.0);
\node[draw, font=\tiny\bfseries, anchor=north west, inner sep=1.5pt] at (8.9,-7.3) {alt};
\node[font=\tiny, anchor=north west] at (9.5,-7.35) {[key already stored]};
\draw[ret] (12.25,-8.0) -- node[ml] {earlier decision (duplicate)} (9.8,-8.0);
\draw[dashed] (8.9,-8.4) -- (13.3,-8.4);
\node[font=\tiny, anchor=north west] at (9.0,-8.45) {[new event]};
\draw[msg] (9.8,-9.1) -- node[ml] {select eligible employee} (12.25,-9.1);
\draw[ret] (12.25,-9.6) -- node[ml] {employee or none} (9.8,-9.6);
\draw[msg] (9.8,-10.3) -- node[ml] {INSERT alert, task, task\_event,\\dispatch\_decision; COMMIT} (12.25,-10.3);
\node[font=\tiny, align=left, anchor=west] at (9.85,-11.3) {no employee: alert and\\decision only, no task};
\draw[ret] (9.8,-12.6) -- node[ml] {DispatchOutcome} (7.35,-12.6);
\draw[ret] (7.35,-13.0) -- node[ml] {DispatchReceipt} (0,-13.0);
\draw[msg] (12.25,-13.4) -- node[ml] {task stream event} (14.7,-13.4);
\end{tikzpicture}
\caption{Sequence diagram: unattended customer to staffing task}
\label{fig:seq-dispatch}
\end{figure}

\subsection{Manager Question}
Figure~\ref{fig:seq-ask} shows UC-04. The web tier never forwards a question without a valid manager session, and the Ask route rejects the question if the manager has already asked ten in the current minute. The orchestrator works only through the scoped registry, which is bound to the manager's store.

\begin{figure}[htbp]
\centering
\begin{tikzpicture}[scale=0.86, transform shape]
\foreach \x/\n/\t in {0/mgr/Manager, 2.45/web/Web tier\\and proxy, 4.9/ask/Ask route, 7.35/orc/Orchestrator, 9.8/agent/Business\\agent, 12.25/tl/Read-only\\tools, 14.7/pg/PostgreSQL} {
  \node[life] (\n) at (\x,0) {\t};
  \draw[dashed, gray] (\n.south) -- (\x,-8.4);
}
\draw[msg] (0,-1.0) -- node[ml] {question} (2.45,-1.0);
\draw[msg] (2.45,-1.6) -- node[ml] {question, interval, filters,\\actor + internal secret} (4.9,-1.6);
\node[draw, font=\tiny, fill=yellow!10, align=left, text width=2.0cm, anchor=north west] at (5.05,-2.05) {manager of this store?\\fewer than 10 per minute?};
\draw[msg] (4.9,-3.25) -- node[ml] {ask(request, scope)} (7.35,-3.25);
\draw[msg] (7.35,-3.7) -- node[ml] {route, then run} (9.8,-3.7);
\draw[msg] (9.8,-4.3) -- node[ml] {invoke(tool, arguments)} (12.25,-4.3);
\draw[msg] (12.25,-4.9) -- node[ml] {read-only query} (14.7,-4.9);
\draw[ret] (14.7,-5.3) -- node[ml] {rows} (12.25,-5.3);
\draw[ret] (12.25,-5.7) -- node[ml] {result with evidence} (9.8,-5.7);
\draw[ret] (9.8,-6.2) -- node[ml] {draft claims} (7.35,-6.2);
\draw[msg] (7.35,-6.6) arc[start angle=90, end angle=-90, radius=0.25] node[font=\tiny, right, pos=0.5] {ground claims};
\draw[ret] (7.35,-7.4) -- node[ml] {answer, citations, limitations} (4.9,-7.4);
\draw[ret] (4.9,-7.8) -- (2.45,-7.8);
\draw[ret] (2.45,-8.2) -- node[ml] {answer shown on dashboard} (0,-8.2);
\end{tikzpicture}
\caption{Sequence diagram: manager asks an analytics question}
\label{fig:seq-ask}
\end{figure}

\subsection{Task Lifecycle}
A dispatched task is created in the \emph{delivered} state. Figure~\ref{fig:task-states} shows the intended lifecycle. The backend does not yet enforce these transitions: it accepts any of the listed states from the task's owner or a manager and records each change as a task event, so enforcing the order is listed as FYP-2 work. The first acknowledgement also stores the acknowledgement time, which System Status reports.

\begin{figure}[htbp]
\centering
\begin{tikzpicture}[font=\small, node distance=8mm and 14mm,
  state/.style={draw, rounded corners=6pt, minimum width=2.4cm, minimum height=8mm, fill=white, align=center}]
\node[circle, fill=black, inner sep=3pt] (init) {};
\node[state, right=10mm of init] (deliv) {delivered};
\node[state, right=of deliv] (ack) {acknowledged\\\scriptsize ack time stored};
\node[state, right=of ack] (prog) {in\_progress};
\node[state, below=14mm of prog] (done) {completed};
\node[state, below=14mm of deliv] (stop) {declined, expired\\or cancelled};
\node[circle, draw, inner sep=3pt, right=10mm of done] (fin) {}; \node[circle, fill=black, inner sep=1.6pt] at (fin) {};
\node[circle, draw, inner sep=3pt, left=10mm of stop] (fin2) {}; \node[circle, fill=black, inner sep=1.6pt] at (fin2) {};
\draw[flow] (init) -- node[lbl, above] {dispatch} (deliv);
\draw[flow] (deliv) -- node[lbl, above] {acknowledge} (ack);
\draw[flow] (ack) -- node[lbl, above] {start} (prog);
\draw[flow] (prog) -- node[lbl, right] {complete} (done);
\draw[flow] (ack) -- node[lbl, sloped, above] {complete} (done);
\draw[flow] (deliv) -- node[lbl, right] {decline} (stop);
\draw[flow] (ack) -- node[lbl, sloped, above] {decline} (stop);
\draw[flow] (done) -- (fin);
\draw[flow] (stop) -- (fin2);
\end{tikzpicture}
\caption{State machine of a staffing task}
\label{fig:task-states}
\end{figure}

\section{Deployment}
Figure~\ref{fig:deployment} shows the deployment the system is built for. The web tier is deployed on Vercel and the database is a hosted PostgreSQL service (Neon). The Python backend and vision pipeline run on one CPU computer near the cameras, because video should not leave the store. At the time of writing the backend has been run only on team laptops; deploying it is part of the FYP-2 plan (Chapter~\ref{ch:conclusion}).

\begin{figure}[htbp]
\centering
\begin{tikzpicture}[font=\small,
  node3d/.style={draw, fill=white, align=center, inner sep=5pt, font=\scriptsize},
  art/.style={draw, fill=white, align=center, font=\scriptsize, inner sep=3pt, text width=3.2cm},
  conn/.style={thin, <->, >={Stealth[length=1.8mm]}},
  cl/.style={font=\tiny, fill=white, inner sep=1pt, align=center}]
% browser
\node[art] (br) at (0,0) {\guillemotleft artifact\guillemotright\\Web pages (React 19)};
\node[art] (nx) at (6.2,0) {\guillemotleft artifact\guillemotright\\Next.js 15 app,\\session and API proxy};
\node[art] (be) at (6.2,-3.6) {\guillemotleft artifact\guillemotright\\FastAPI backend (Uvicorn):\\API, agents, dispatch};
\node[art] (cv) at (6.2,-5.5) {\guillemotleft artifact\guillemotright\\Vision pipeline jobs:\\YOLOv8n, best\_320, pose};
\node[art] (cam) at (0,-4.55) {\guillemotleft device\guillemotright\\CCTV cameras or\\recorded video};
\node[art] (pg) at (12.2,-3.6) {\guillemotleft database\guillemotright\\PostgreSQL 17};
\begin{scope}[on background layer]
\node[node3d, fill=gray!8, fit=(br), inner sep=8pt, label={[font=\scriptsize]above:\guillemotleft device\guillemotright\ Browser}] (dev) {};
\node[node3d, fill=gray!8, fit=(nx), inner sep=8pt, label={[font=\scriptsize]above:\guillemotleft cloud\guillemotright\ Vercel}] (vc) {};
\node[node3d, fill=gray!8, fit=(be)(cv), inner sep=8pt, label={[font=\scriptsize]below:\guillemotleft device\guillemotright\ Store application host (CPU only)}] (host) {};
\node[node3d, fill=gray!8, fit=(pg), inner sep=8pt, label={[font=\scriptsize]above:\guillemotleft cloud\guillemotright\ Neon}] (neon) {};
\end{scope}
% external
\node[art] (ext) at (12.2,0) {\guillemotleft service\guillemotright\\Hugging Face router,\\ElevenLabs (optional)};
\draw[conn] (dev) -- node[cl, above] {HTTPS} (vc);
\draw[conn] (vc) -- node[cl, left, pos=0.35] {HTTPS, internal secret,\\actor headers} (vc |- host.north);
\draw[conn] (host.east |- pg) -- node[cl, above] {TLS (sslmode=require)} (neon);
\draw[conn] (host.north east) -- node[cl, sloped, above] {HTTPS} (ext);
\draw[conn] (cam) -- node[cl, above] {stream / files} (host.west |- cam);
\end{tikzpicture}
\caption{Deployment diagram}
\label{fig:deployment}
\end{figure}


